% Generated draft: requires mathematical and source review.
\documentclass[11pt]{article}
\usepackage[margin=25mm]{geometry}
\usepackage{fontspec}
\setmainfont{Latin Modern Roman}
\usepackage{amsmath,amssymb,mathtools,unicode-math}
\setmathfont{Latin Modern Math}
\usepackage{enumitem,needspace,xurl,hyperref}
\setlength{\parindent}{0pt}
\setlength{\parskip}{5pt}
\setlength{\emergencystretch}{3em}
\newcommand{\N}{\mathbb{N}}
\newcommand{\Z}{\mathbb{Z}}
\newcommand{\Q}{\mathbb{Q}}
\newcommand{\R}{\mathbb{R}}
\newcommand{\C}{\mathbb{C}}
\newcommand{\F}{\mathbb{F}}
\newcommand{\A}{\mathbb{A}}
\newcommand{\PP}{\mathbb{P}}
\newcommand{\E}{\mathbb{E}}
\newcommand{\eps}{\varepsilon}
\newcommand{\dd}{\,\mathrm{d}}
\begin{document}
\section*{UnsolvedMath ID: 30000231}
\textbf{Source problem code:} catalogue entry
\subsection{Definitions and mathematical statement}
Let \((\mathcal M,g)\) be the Schwarzschild or Reissner--Nordstr\"om exterior intended in the source, with area-radius function \(r\) and outgoing and ingoing null coordinates \(u,v\). Let \(\phi\) denote the zeroth spherical harmonic of a solution of the linear equation \(\Box_g\phi=0\) arising from compactly supported Cauchy data. Write
\[
\Phi_{\mathcal I^+}(u)=\lim_{v\to\infty}r(u,v)\phi(u,v)
\]
for its radiation field where that limit is defined. The source supplies the late-time bounds
\[
|\Phi_{\mathcal I^+}(u)|\leq Cu^{-2},\qquad
|\partial_u\Phi_{\mathcal I^+}(u)|\leq C_\varepsilon u^{-3+\varepsilon},
\]
and, in the exterior region \(r(u,v)\leq R\),
\[
|\phi(u,v)|\leq C_{R,\varepsilon}v^{-3+\varepsilon},\qquad
|\partial_v\phi(u,v)|\leq C_{R,\varepsilon}v^{-3+\varepsilon}.
\]
The conventional reading is for each \(\varepsilon>0\), each fixed admissible finite \(R\), and sufficiently large positive times, with constants depending on the data and indicated parameters. The exact coordinate normalization, background parameters, and regularity hypotheses must be recovered from the report. In particular, subextremal and extremal Reissner--Nordstr\"om backgrounds must not be identified without checking the applicable estimates.

The broad question is whether these decay rates are sufficient to control an intended nonlinearity by perturbation of the linear problem. The explicit model asks for the small-data behavior of
\[
\Box_g\phi=|\phi|^p,\qquad 1+\sqrt2\leq p\leq4.
\]
Both endpoints are included. The source calls this the remaining range in its historical context; this is not an assertion that the same range remains unresolved now. It does not supply the Cauchy-data norm, exact global-existence, blow-up, lifespan, or scattering target, or a class of general nonlinear perturbations. The sign convention for \(\Box_g\) is also not reproduced.

The displayed linear estimates concern the zeroth spherical harmonic. For a general nonlinear field, projecting \(|\phi|^p\) onto that harmonic does not simply replace \(\phi\) by its projection; a spherical-symmetry restriction or higher-mode estimates would need to be specified. This draft does not infer an all-mode nonlinear theorem solely from the quoted bounds.

The supplied literature assessment checked 2026-08-21 reports later global small-data and critical-behavior results in substantial ranges for these models, while saying that the broad question for general nonlinear perturbations is not fully settled. Its precise parameter, symmetry, and exponent hypotheses are not supplied.
\subsection{Short English statement}
Can the quoted linear decay bounds on black-hole exteriors support perturbative treatment of nonlinear waves, especially the power model for \(1+\sqrt2\leq p\leq4\)? The source's historical range, harmonic restriction, background parameters, and small-data notion need to be kept distinct.
\subsection{Sources}
\begin{itemize}
\item Local UnsolvedMath registry, permanent ID 30000231, catalogue entry, \emph{Decay Rates for Nonlinear Waves on Black-Hole Backgrounds}.
\item \emph{Nonlinear Evolution Problems} (2006), Oberwolfach Reports, pp. 1371--1430, DOI \url{https://doi.org/10.4171/owr/2005/25}, estimates (3)--(4) as identified in the original wording.
\item Dafermos--Rodnianski, supplied label ``nonlinear waves on black-hole backgrounds'': \url{https://arxiv.org/abs/math/0503024}.
\item Lai--Zhou, supplied label ``nonlinear waves in Schwarzschild spacetime'': \url{https://doi.org/10.1016/j.matpur.2023.02.009}.
\item Supplied source-verification claim restores the bounds, model, and historical exponent range; supplied literature triage checked 2026-08-21.
\end{itemize}
\end{document}
